\clearpage
\thispagestyle{empty}
\SourceLeaf{329}\pdfbookmark[0]{On Parabolic Motion}{parabolische-bewegung}
\vspace*{1.6in}
\begin{center}
{\Large ON PARABOLIC MOTION.\par}
\end{center}
\clearpage
\fancyhead[L]{\small ON PARABOLIC MOTION}
\SourceLeaf{331}\pdfbookmark[1]{Posthumous Papers and Correspondence}{parabolisch-nachlass}
\begin{center}
{\large POSTHUMOUS PAPERS AND CORRESPONDENCE.\par}
\bigskip
[I. On the Determination of Cometary Orbits.]\par
\medskip
\pdfbookmark[2]{1. Parabolic Motion. Time from Two Radius Vectors and the Chord.}{parabolisch-001}[1.]\par
\medskip
Parabolic motion.\par
\medskip

Time from two radius vectors and the chord.
\medskip

[Notation as in Theoria motus corporum coelestium, Article~108.]
\end{center}
\[
\begin{aligned}
r&=\tfrac{1}{2}p(1+\theta\theta),\qquad r'=\tfrac{1}{2}p(1+\theta'\theta')\\
\rho&=p(\theta'-\theta)\eta,\qquad \eta=\tfrac{1}{2}\BareRoot{(4+(\theta'+\theta)^2)}
\end{aligned}
\]
\[
r'+r=\tfrac{1}{4}p(4+2\theta\theta+2\theta'\theta')=\tfrac{1}{4}p(4\eta\eta+(\theta'-\theta)^2)
\]
\[
r'+r+\rho=\tfrac{1}{4}p(2\eta+\theta'-\theta)^2,\qquad
r'+r-\rho=\tfrac{1}{4}p(2\eta-\theta'+\theta)^2.
\]

If, then, we set
\[
\BareRoot{(r'+r+\rho)}=m,\qquad \pm\BareRoot{(r'+r-\rho)}=n,
\]
taking the upper or lower sign according as the motion amounts to less or more than \(180^0\), we have
\[
\begin{aligned}
m&=\BareRoot{p}\,.(\eta+\tfrac{1}{2}\theta'-\tfrac{1}{2}\theta)\\
&=\frac{\rho}{(\theta'-\theta)\BareRoot{p}}+\frac{(\theta'-\theta)\BareRoot{p}}{2}\\[.5ex]
n&=\BareRoot{p}\,.(\eta-\tfrac{1}{2}\theta'+\tfrac{1}{2}\theta)\\
&=\frac{\rho}{(\theta'-\theta)\BareRoot{p}}-\frac{(\theta'-\theta)\BareRoot{p}}{2}\mathpunct{\cdot}
\end{aligned}
\]
Hence
\[
m-n=(\theta'-\theta)\BareRoot{p};
\]
furthermore,
\[
r'-r=\tfrac{1}{2}p(\theta'-\theta)(\theta'+\theta),
\]
so that
\[
\frac{2(r'-r)}{m-n}=(\theta'+\theta)\BareRoot{p},
\]
\SourceLeaf{332}[which gives]
\[
\begin{aligned}
\theta\BareRoot{p}&=\frac{2(r'-r)-(m-n)^2}{2(m-n)}\\[.5ex]
\theta'\BareRoot{p}&=\frac{2(r'-r)+(m-n)^2}{2(m-n)}\mathpunct{\cdot}
\end{aligned}
\]

We also have
\[
\begin{aligned}
\rho\rho&=pp(\theta'-\theta)^2\eta\eta\\
&=pp(\theta'-\theta)^2+\tfrac{1}{4}pp(\theta'-\theta)^2(\theta'+\theta)^2\\
&=pp(\theta'-\theta)^2+(r'-r)^2
\end{aligned}
\]
[or, since \(2\rho=mm-nn\),]
\[
\frac{\rho\rho-(r'-r)^2}{(m-n)^2}=p=\frac{(mm-nn)^2-4(r'-r)^2}{4(m-n)^2}\mathpunct{\cdot}
\]

[If \(t\) and \(t'\) are the respective times elapsed since the passage through perihelion, we thus have]
\[
\begin{aligned}
t&=\frac{p^{\frac{3}{2}}}{6k}(3\theta+\theta^3)=\frac{p}{2k}\theta\BareRoot{p}+\frac{1}{6k}(\theta\BareRoot{p})^3\\[.5ex]
&=\frac{1}{48k(m-n)^3}\Bigl\{(2(r'-r)-(m-n)^2)^3\\
&\qquad-3(2(r'-r)-(m-n)^2)(4(r'-r)^2-(mm-nn)^2)\Bigr\}\\[.5ex]
&=\frac{1}{12k(m-n)^3}\Bigl\{-4(r'-r)^3+3(r'-r)(m-n)^2(mm+nn)\\
&\qquad-(m-n)^4(mm+mn+nn)\Bigr\}
\end{aligned}
\]
[and similarly]
\[
\begin{aligned}
t'&=\frac{1}{12k(m-n)^3}\Bigl\{-4(r'-r)^3+3(r'-r)(m-n)^2(mm+nn)\\
&\qquad+(m-n)^4(mm+mn+nn)\Bigr\}.
\end{aligned}
\]

Hence
\[
\begin{aligned}
\mathrm{[I]}\quad&t'-t=\frac{1}{6k}(m^3-n^3)\\[1ex]
\mathrm{[II]}\quad&\tfrac{1}{2}(t'+t)=\frac{1}{6k}\left\{-2\left(\frac{r'-r}{m-n}\right)^2+\frac{3(r'r'-rr)}{m-n}\right\}\mathpunct{\cdot}
\end{aligned}
\]

The latter formula alone is more easily obtained as follows:
\[
\begin{aligned}
&\tfrac{1}{2}(t'+t)=\frac{p^{\frac{3}{2}}}{12k}(\theta'+\theta)(3+\theta'\theta'-\theta'\theta+\theta\theta)\\
&\qquad=\frac{p^{\frac{3}{2}}}{12k}(\theta'+\theta)\left(\frac{3(r'+r)}{p}-\tfrac{1}{2}(\theta'+\theta)^2\right)\\[.5ex]
&\qquad=\frac{r'+r}{4k}\cdot\frac{2(r'-r)}{m-n}-\frac{1}{24k}\left(\frac{2(r'-r)}{m-n}\right)^3\\[.5ex]
\mathrm{[II^{*}]}\quad&\qquad=\frac{1}{2k}\cdot\frac{r'r'-rr}{m-n}-\frac{1}{3k}\left(\frac{r'-r}{m-n}\right)^3\mathpunct{\cdot}
\end{aligned}
\]
\begin{center}\rule{.3\linewidth}{.4pt}\end{center}
